\documentclass[11pt]{article}
\usepackage[margin=1in]{geometry}
\usepackage{amsmath,amssymb,amsthm}
\usepackage[hidelinks]{hyperref}

\newtheorem{theorem}{Theorem}
\newtheorem{lemma}{Lemma}
\newtheorem{corollary}{Corollary}
\newtheorem{remark}{Remark}

\newcommand{\weak}{\sigma}
\newcommand{\cc}{\mathfrak c}
\newcommand{\ellone}{\ell_1}

\title{Positive families for the Schur weak-\(\aleph\) question}
\author{FA/Banach run packet for arXiv:1412.1748}
\date{2026-07-19}

\begin{document}
\maketitle

\section*{Source question}

The source paper is S. Gabriyelyan, J. Kakol, W. Kubis and W. Marciszewski,
\emph{Networks for the weak topology of Banach and Frechet spaces},
arXiv:1412.1748.
After proving that \(\ell_1(\Gamma)\) is an \(\aleph\)-space in the weak
topology if and only if \(|\Gamma|\leq\mathfrak c\), the authors ask:

\begin{quote}
Let \(E\) be a Banach space with the Schur property and satisfy (i)--(iii)
of Proposition 3.4. Is \(E\) an \(\aleph\)-space in the weak topology?
\end{quote}

This packet does not settle the question in full.  It records two positive
families which are stable and useful test cases: weakly compactly generated
Schur spaces, and Banach spaces isomorphic to subspaces of \(\ell_1(\Gamma)\)
for \(|\Gamma|\leq\mathfrak c\).

\section*{Idea of the proof}

The source already proves the hard nonseparable model case
\(\ell_1(\Gamma)\), \(|\Gamma|\leq\mathfrak c\).  The first observation is
that the weak topology on a Banach subspace is exactly the subspace topology
inherited from the ambient weak topology, because Hahn--Banach extends every
functional on the subspace.  Since \(\aleph\)-spaces are hereditary to
subspaces, the whole subspace family follows immediately.

The WCG family is different.  If \(E\) has the Schur property, then weakly
compact subsets of \(E\) are norm compact.  A weakly compact generating set
therefore has separable norm-closed span, so a WCG Schur space is separable.
For a separable Schur space, a countable norm base is a countable \(k\)-network
for the weak topology: weakly compact sets are norm compact and weakly open
sets are norm open.

\begin{lemma}\label{lem:subspace}
Let \(Y\) be a closed Banach subspace of a Banach space \(X\).  Then the weak
topology \(\sigma(Y,Y')\) is the topology inherited by \(Y\) from
\((X,\sigma(X,X'))\).
\end{lemma}

\begin{proof}
The inherited topology is generated by the restrictions \(x'|_Y\), \(x'\in X'\).
It is therefore no stronger than \(\sigma(Y,Y')\).  Conversely, by the
Hahn--Banach theorem, every \(y'\in Y'\) is the restriction of some
\(x'\in X'\).  Hence the two initial topologies are generated by the same
linear functionals on \(Y\).
\end{proof}

\begin{lemma}\label{lem:hereditary}
Every subspace of an \(\aleph\)-space is an \(\aleph\)-space.
\end{lemma}

\begin{proof}
Let \(\mathcal N=\bigcup_n\mathcal N_n\) be a \(\sigma\)-locally finite
\(k\)-network in \(X\), and let \(Y\subset X\).  The family
\(\{N\cap Y:N\in\mathcal N\}\) is still \(\sigma\)-locally finite in \(Y\).
If \(K\subset U\) with \(K\) compact in \(Y\) and \(U\) open in \(Y\), write
\(U=V\cap Y\) with \(V\) open in \(X\).  Then \(K\) is compact in \(X\), so
there are \(N_1,\dots,N_m\in\mathcal N\) such that
\[
  K\subset \bigcup_{j=1}^m N_j\subset V.
\]
Intersecting with \(Y\) gives the required finite \(k\)-network refinement in
\(Y\).
\end{proof}

\begin{theorem}\label{thm:l1subspaces}
Let \(\Gamma\) be a set with \(|\Gamma|\leq\mathfrak c\).  If a Banach space
\(E\) is isomorphic to a closed subspace of \(\ell_1(\Gamma)\), then
\((E,\sigma(E,E'))\) is an \(\aleph\)-space.
\end{theorem}

\begin{proof}
By Theorem 1.2 of the source paper, \(\ell_1(\Gamma)\) endowed with its weak
topology is an \(\aleph\)-space whenever \(|\Gamma|\leq\mathfrak c\).  If
\(F\subset\ell_1(\Gamma)\) is a closed subspace, Lemma~\ref{lem:subspace}
identifies \((F,\sigma(F,F'))\) with a subspace of
\((\ell_1(\Gamma),\sigma(\ell_1(\Gamma),\ell_1(\Gamma)'))\).  Lemma
\ref{lem:hereditary} gives that \(F\) is weakly an \(\aleph\)-space.  Finally,
a Banach-space isomorphism \(E\to F\) is a homeomorphism for the weak
topologies, since its adjoint identifies the corresponding dual functionals.
\end{proof}

\begin{lemma}\label{lem:separable-schur}
Every separable Banach space with the Schur property is weakly
\(\aleph_0\).
\end{lemma}

\begin{proof}
Let \(\mathcal B\) be a countable base for the norm topology of \(E\).  We
claim that \(\mathcal B\) is a \(k\)-network for the weak topology.
Let \(K\subset U\), where \(K\) is weakly compact and \(U\) is weakly open.
By the Eberlein--Smulian theorem and the Schur property, every sequence in
\(K\) has a norm-convergent subsequence; hence \(K\) is norm compact.  Also
\(U\) is norm open.  A finite subfamily of the norm base covers \(K\) and is
contained in \(U\).  Thus \(\mathcal B\) is a countable \(k\)-network for
\((E,\sigma(E,E'))\).
\end{proof}

\begin{corollary}\label{cor:wcg}
Every weakly compactly generated Banach space with the Schur property is
weakly \(\aleph_0\).  In particular, Question 4.9 has an affirmative answer
for WCG Schur spaces satisfying the source paper's conditions (i)--(iii).
\end{corollary}

\begin{proof}
Let \(K\) be a weakly compact subset whose closed linear span is \(E\).  By
the Schur property and Eberlein--Smulian, \(K\) is norm compact, hence norm
separable.  The closed linear span of a separable set is separable, so \(E\)
is separable.  Lemma~\ref{lem:separable-schur} applies.
\end{proof}

\section*{Limitations and obstruction}

The result above is a genuine positive subcase but not a full solution of
Question 4.9.  The source proof for \(\ell_1(\Gamma)\) uses a special annulus
localization lemma: near every point of a thin norm annulus one can find a
weak neighborhood whose intersection with that annulus has small norm
diameter.  This follows from the coordinate structure of \(\ell_1(\Gamma)\).
The Schur property alone only controls weakly convergent sequences and weakly
compact sets; it does not by itself give such local norm control on arbitrary
bounded annuli.  Thus the quotient/general Schur-space case remains open in
this packet.

\section*{Novelty check}

Local run indexes were searched for arXiv:1412.1748, the source title,
``Schur property'', ``weakly aleph'', ``countable Tychonoff'', and
``\(\ell_1(\Gamma)\) weak topology''.  Web searches for the exact Schur
question and for later occurrences of ``Banach space with the Schur property''
and ``weakly aleph'' found the source paper but no exact later answer to
Question 4.9.  A later paper by Ferrando--Gabriyelyan--Kakol titled
\emph{Functional Characterizations of Countable Tychonoff Spaces} has adjacent
metadata about function-space characterizations, but the public abstract and
local corpus did not expose an exact answer to this Schur-space question.

\section*{References}

\begin{thebibliography}{9}
\bibitem{source}
S. Gabriyelyan, J. Kakol, W. Kubis and W. Marciszewski,
\emph{Networks for the weak topology of Banach and Frechet spaces},
arXiv:1412.1748.

\bibitem{fabian}
M. Fabian, P. Habala, P. Hajek, V. Montesinos, J. Pelant and V. Zizler,
\emph{Functional Analysis and Infinite-Dimensional Geometry},
CMS Books in Mathematics, Springer, 2001.
\end{thebibliography}

\end{document}
